
\exercise
Suppose $T \in \L(V\1, W)$. Show that $T = 0$ if and only if all singular values of $T$ are $0$.


\exercise
Suppose $T \in \L(V\1, W)$ and $s > 0$. Prove that $s$ is a singular value of $T$ if and only if there exist nonzero vectors $v \in V$ and
$w \in W$ such that \label{2Schmidt}
\[
Tv = sw \andd T^* w = sv.
\]
\exercisedisplayend
\exercisecomment{The vectors $v, w$ satisfying both equations above are called a \marbf{Schmidt pair}. Erhard Schmidt introduced the concept of singular values in 1907.\index{Schmidt pair}\index{Schmidt, Erhard}}


\exercise
Give an example of $T \in \L\paren[\big]{\C2}$ such that $0$ is the only eigenvalue of $T$ and the singular values of $T$ are $5, 0$.


\exercise
Suppose that $T \in \L(V\1, W)$, $s_1$ is the largest singular value of $T$, and $s_n$ is the smallest singular value of $T\3$. Prove that\label{3interval}
\[
\br[\big]{ \norm{Tv} : v \in V \text{ and } \norm{v} = 1 } = [s_n, s_1].
\]
\exercisedisplayend


\exercise
Suppose $T \in \L\paren[\big]{\C{2}}$ is defined by $T(x, y) = (-4y, x)$. Find the singular values of $T\3$.


\exercise
Find the singular values of the differentiation operator $D \in \L\bigl(\P\1_2(\R{})\bigr)$ defined by $Dp = p'\!$, where the inner product on $\P\1_2(\R{})$ is as in Example \ref{obP2R}.


\exercise
Suppose that $T \in \L(V)$ is self-adjoint or that $\F{}= \C{}$ and $T \in \L(V)$ is normal. Let $\la_1, \ldots, \la_n$ be the eigenvalues of $T$, each included in this list as many times as the dimension of the corresponding eigenspace. Show that the singular values of $T$ are
$\abs{\la_1}, \ldots, \abs{\la_n}$, after these numbers have been sorted into decreasing order.\label{7singsp}


\exercise
Suppose $T \in \L(V\1, W)$. Suppose $s_1 \ge  s_2 \ge \cdots \ge s_m >0$ and
$e_1, \dots, e_m$ is an orthonormal list in $V$ and $f_1, \dots, f_m$ is an orthonormal list in~$W$ such that \label{4SVD}
\[
Tv = s_1 \ip{v, e_1} f_1 + \dots + s_m \ip{v, e_m} f_m
\]
for every $v \in V\1$.
\begin{subtheorem}
\item
Prove that $f_1, \dots, f_m$ is an orthonormal basis of $\ran T\3$.
\item
Prove that $e_1, \dots, e_m$  is an orthonormal basis of $(\ker T)^\perp$.
\item
Prove that $s_1, \ldots, s_m$ are the positive singular values of $T\3$.
\item
Prove that if $k \in \br{1, \ldots, m}$, then $e_k$ is an eigenvector of $T^*T$ with corresponding eigenvalue ${s_k}^{\2}\1$.
\item
Prove that
\[
TT^*w = {s_1}^{\2} \ip{w, f_1} f_1 + \cdots + {s_m}^{\2} \ip{w, f_m} f_m
\]
for all $w \in W\3$.
\end{subtheorem}
\exercisesubtheoremend


\exercise
Suppose $T \in \L(V\1, W)$. Show that $T$ and $T^*$ have the same positive singular values. \label{4TstarSing}


\exercise
Suppose $T \in \L(V\1, W)$ has singular values $s_1, \ldots, s_n$. Prove that if $T$ is an  invertible linear map, then $T^{-1}$ has singular values\label{7singinv}
\[
\frac{1}{s_n}, \ldots, \frac{1}{s_1}.
\]
\exercisedisplayend


\exercise
Suppose that $T \in \L(V\1, W)$ and $v_1, \ldots, v_n$ is an orthonormal basis of $V\1$. Let $s_1, \ldots, s_n$ denote the singular values of $T\3$.\label{7sums2}
\begin{subtheorem}
\item
Prove that
$\norm{Tv_1}^2 + \cdots + \norm{Tv_n}^2 = {s_1}^{\2} + \cdots + {s_n}^{\2}\1$.
\item
Prove that if $W = V$ and $T$ is a positive operator, then
\[
\ip{Tv_1, v_1} + \cdots + \ip{Tv_n, v_n} = s_1 + \cdots + s_n.
\]
\exercisedisplayend
\end{subtheorem}
\exercisecomment{See the comment after Exercise \ref{6sumT2} in Section \ref{7001}.}


\exercise
\begin{SUBtheorem}
\item
Give an example of a finite-dimensional vector space and an operator $T$ on it
such that the singular values of $T^2$ do not equal the squares of the singular values of~$T\3$.
\item
Suppose $T \in \L(V)$ is normal. Prove that the singular values of $T^2$ equal the squares of the singular values of~$T\3$.
\end{SUBtheorem}


\exercise
Suppose $T\3_1, T\3_2 \in \L(V)$.  Prove that $T\3_1$ and $T\3_2$ have the
same singular values if and only if there exist unitary operators $S_1, S_2
\in \L(V)$ such that $T\3_1 = S_1 T\3_2 S_2$.


\exercise
Suppose $T \in \L(V\1,W)$.  Let $s_n$ denote the smallest singular value of~$T\3$.
Prove that $s_n \| v \| \le \| Tv \|$ for every $v \in V\1$. \label{3smallSV}


\exercise
Suppose $T \in \L(V)$ and $s_1 \ge \cdots \ge s_n$ are the singular values of $T\3$.
Prove that if $\la$ is an eigenvalue of $T$, then $s_1  \ge |\la| \ge s_n$.


\exercise
Suppose $T \in \L(V\1, W)$. Prove that $\paren[\big]{T^*}^\dagger = \paren[\big]{T^\dagger}^*\1$.\label{7Tdagstar}\index{pseudoinverse|(}
\exercisecomment{Compare the result in this exercise to the analogous result for invertible linear maps \textup{\big[}see \ref{7propadjoint}\uppar{\,f\,}\textup{\big]}.}


\exercise
Suppose $T \in \L(V)$. Prove that $T$ is self-adjoint if and only if $T^\dagger$ is self-adjoint.\index{pseudoinverse|)}


